\chapter{Linear Algebra and Optimization}
\label{ch:linear-algebra}

\epigraph{``Statistics is the science of data. Linear algebra is the language of statistics.''}{}

\learninggoal{
	\item Review the essential linear algebra: vectors, matrices, eigenvalues, SVD.
	\item Understand gradients, convexity, and gradient-based optimization.
	\item Derive the normal equations for linear regression.
	\item Connect matrix factorizations to dimensionality reduction.
}

\section{Vectors and Matrices}

A vector $\mathbf{x} \in \R^d$ is an ordered list of $d$ numbers. A matrix $\mathbf{A} \in \R^{m \times n}$ is a rectangular array.

\subsection{Key Operations}

\begin{itemize}
	\item \textbf{Inner product:} $\mathbf{x}^\top \mathbf{y} = \sum_{i=1}^d x_i y_i$
	\item \textbf{Norm:} $\|\mathbf{x}\|_p = (\sum_i |x_i|^p)^{1/p}$, with $\|\mathbf{x}\|_2 = \sqrt{\mathbf{x}^\top \mathbf{x}}$
	\item \textbf{Matrix--vector product:} $(\mathbf{A}\mathbf{x})_i = \sum_j A_{ij} x_j$
	\item \textbf{Matrix--matrix product:} $(\mathbf{A}\mathbf{B})_{ij} = \sum_k A_{ik} B_{kj}$
	\item \textbf{Transpose:} $(\mathbf{A}^\top)_{ij} = A_{ji}$
	\item \textbf{Trace:} $\operatorname{tr}(\mathbf{A}) = \sum_i A_{ii}$
\end{itemize}

\subsection{Special Matrices}

\begin{longtable}{p{3cm}p{5cm}p{4cm}}
	\toprule
	\textbf{Type}     & \textbf{Definition}                                                 & \textbf{Property}        \\
	\midrule
	Symmetric         & $\mathbf{A} = \mathbf{A}^\top$                                      & Real eigenvalues         \\
	Positive definite & $\mathbf{x}^\top \mathbf{A} \mathbf{x} > 0$ for $\mathbf{x} \neq 0$ & All eigenvalues positive \\
	Orthogonal        & $\mathbf{Q}^\top \mathbf{Q} = \mathbf{I}$                           & Preserves norms          \\
	Diagonal          & $A_{ij} = 0$ for $i \neq j$                                         & Easy to invert           \\
	\bottomrule
\end{longtable}

\section{Eigenvalues and Eigenvectors}

\begin{definition}[Eigenvalue Decomposition]
	For a square matrix $\mathbf{A} \in \R^{d \times d}$, $\mathbf{v}$ is an \textbf{eigenvector} with \textbf{eigenvalue} $\lambda$ if:
	\[
		\mathbf{A} \mathbf{v} = \lambda \mathbf{v}
	\]
	Equivalently: $\mathbf{A} \mathbf{V} = \mathbf{V} \boldsymbol{\Lambda}$, where $\mathbf{V}$'s columns are eigenvectors and $\boldsymbol{\Lambda}$ is diagonal with eigenvalues.
\end{definition}

Key facts:
\begin{itemize}
	\item $\operatorname{tr}(\mathbf{A}) = \sum_i \lambda_i$
	\item $\det(\mathbf{A}) = \prod_i \lambda_i$
	\item A symmetric matrix can be diagonalized as $\mathbf{A} = \mathbf{Q} \boldsymbol{\Lambda} \mathbf{Q}^\top$ with orthogonal $\mathbf{Q}$.
	\item The eigenvalues of $\mathbf{A}^{-1}$ are $1/\lambda_i$.
\end{itemize}

Eigenvalues are central to PCA (Chapter~\ref{ch:dimensionality-reduction}), where the eigenvectors of the covariance matrix give the principal components.

\section{Singular Value Decomposition}

The SVD generalizes eigenvalue decomposition to non-square matrices.

\begin{definition}[SVD]
	Any matrix $\mathbf{A} \in \R^{m \times n}$ can be factored as:
	\[
		\mathbf{A} = \mathbf{U} \boldsymbol{\Sigma} \mathbf{V}^\top
	\]
	where $\mathbf{U} \in \R^{m \times m}$ and $\mathbf{V} \in \R^{n \times n}$ are orthogonal, and $\boldsymbol{\Sigma} \in \R^{m \times n}$ is diagonal with $\sigma_1 \geq \sigma_2 \geq \cdots \geq \sigma_{\min(m,n)} \geq 0$.
\end{definition}

The singular values $\sigma_i$ and the SVD are the mathematical foundation for:

\begin{itemize}
	\item \textbf{PCA:} the right singular vectors $\mathbf{V}$ are the principal components.
	\item \textbf{Matrix approximation:} keeping only the top $k$ singular values gives the best rank-$k$ approximation (Eckart--Young theorem).
	\item \textbf{Pseudoinverse:} $\mathbf{A}^+ = \mathbf{V} \boldsymbol{\Sigma}^+ \mathbf{U}^\top$ solves least squares.
	\item \textbf{Rank and condition number:} $\text{rank}(\mathbf{A}) = \# \text{ nonzero } \sigma_i$, $\kappa(\mathbf{A}) = \sigma_{\max} / \sigma_{\min}$.
\end{itemize}

\section{Calculus and Gradients}

\subsection{Derivatives of Vector Functions}

The gradient of a scalar function $f: \R^d \to \R$ is:

\[
	\nabla f(\mathbf{x}) = \begin{pmatrix}
		\frac{\partial f}{\partial x_1} & \cdots & \frac{\partial f}{\partial x_d}
	\end{pmatrix}^\top
\]

The Hessian $\nabla^2 f(\mathbf{x})$ is the $d \times d$ matrix of second derivatives.

\subsection{Useful Matrix Derivatives}

\begin{itemize}
	\item $\nabla_{\mathbf{x}} (\mathbf{a}^\top \mathbf{x}) = \mathbf{a}$
	\item $\nabla_{\mathbf{x}} (\mathbf{x}^\top \mathbf{A} \mathbf{x}) = (\mathbf{A} + \mathbf{A}^\top) \mathbf{x}$
	\item $\nabla_{\mathbf{x}} \|\mathbf{A}\mathbf{x} - \mathbf{b}\|_2^2 = 2\mathbf{A}^\top (\mathbf{A}\mathbf{x} - \mathbf{b})$
	\item $\nabla_{\mathbf{X}} \operatorname{tr}(\mathbf{A}\mathbf{X}) = \mathbf{A}^\top$
\end{itemize}

\begin{example}[Deriving the Normal Equations]
	Minimize $L(\mathbf{w}) = \|\mathbf{X}\mathbf{w} - \mathbf{y}\|_2^2$:
	\[
		\nabla L = 2\mathbf{X}^\top (\mathbf{X}\mathbf{w} - \mathbf{y}) = 0
	\]
	\[
		\mathbf{X}^\top \mathbf{X} \mathbf{w} = \mathbf{X}^\top \mathbf{y}
	\]
	\[
		\hat{\mathbf{w}} = (\mathbf{X}^\top \mathbf{X})^{-1} \mathbf{X}^\top \mathbf{y}
	\]
\end{example}

\section{Convexity}

\begin{definition}[Convex Function]
	A function $f: \R^d \to \R$ is \textbf{convex} if for all $\mathbf{x}, \mathbf{y}$ and $\lambda \in [0,1]$:
	\[
		f(\lambda \mathbf{x} + (1-\lambda) \mathbf{y}) \leq \lambda f(\mathbf{x}) + (1-\lambda) f(\mathbf{y})
	\]
	Equivalently (if twice differentiable): $\nabla^2 f(\mathbf{x}) \succeq 0$ (positive semidefinite Hessian).
\end{definition}

Convexity guarantees that any local minimum is a global minimum. Linear regression, logistic regression, and SVM (with convex loss) all optimize convex functions. Neural networks are non-convex.

\section{Gradient Descent}

When the normal equations are not available (e.g., logistic regression, neural networks), we use iterative optimization.

\[
	\mathbf{w}^{(t+1)} = \mathbf{w}^{(t)} - \eta \nabla L(\mathbf{w}^{(t)})
\]

\subsection{Learning Rate}

The learning rate $\eta$ controls step size:

\begin{itemize}
	\item Too large: divergence.
	\item Too small: slow convergence.
	\item Adaptive methods (Adam, RMSprop) adjust $\eta$ per parameter.
\end{itemize}

\subsection{Stochastic Gradient Descent}

Instead of computing the gradient on all $n$ data points (which is $O(nd)$ per step), SGD approximates it with a single data point:

\[
	\nabla L(\mathbf{w}) \approx \nabla \ell(y_i, f(\mathbf{x}_i; \mathbf{w}))
\]

Mini-batch SGD uses $B$ points (e.g., $B = 32$), balancing noise and computation.

\begin{principle}[SGD Convergence]
	SGD converges to a global minimum for convex losses and to a good local minimum for non-convex losses, provided the learning rate decays appropriately: $\sum_t \eta_t = \infty$ and $\sum_t \eta_t^2 < \infty$.
\end{principle}

\section{The Chain Rule and Backpropagation}

For nested functions $L = \ell(f(g(\mathbf{x}; \mathbf{w})))$:

\[
	\frac{\partial L}{\partial w_{ij}} = \frac{\partial L}{\partial f_k} \cdot \frac{\partial f_k}{\partial g_j} \cdot \frac{\partial g_j}{\partial w_{ij}}
\]

This is the chain rule, applied systematically through a computation graph. Backpropagation is simply the chain rule with caching of intermediate results.

\section{Exercises}

\begin{exercise}
	Compute the eigenvalues and eigenvectors of $\mathbf{A} = \begin{pmatrix} 2 & 1 \\ 1 & 2 \end{pmatrix}$.
\end{exercise}

\begin{exercise}
	Prove that $\|\mathbf{A}\mathbf{x}\|_2 \leq \sigma_{\max}(\mathbf{A}) \|\mathbf{x}\|_2$ for any vector $\mathbf{x}$.
\end{exercise}

\begin{exercise}
	Take one step of gradient descent on $f(x) = x^4 - 3x^3 + 2$ starting at $x = 2$ with $\eta = 0.01$. What is the new value of $x$?
\end{exercise}

\begin{exercise}
	Show that the sum of convex functions is convex. Use this to explain why the empirical risk for logistic regression is convex.
\end{exercise}
